% SPDX-License-Identifier: CC-BY-4.0
% Copyright (C) 2026 the Langlandscape team
% Stage 00 (introduction): how the route makes JL computable, and how we implemented it in HACS, for non-specialists.
% Written from our sources only; no new mathematics, no status upgrades. Every number is computed (checked where it equals an
% independent reference); nothing is certified. Sources: plan/hacs/HACS-JL-PLAN.md, plan/hacs/REPORT.md, the lane closures
% plan/hacs/outbox/20261115T0137Z_laneA-closure.md, 20261115T0150Z_laneB-closure.md (with its errata),
% 20261115T0830Z_laneC-summary.md, 20261115T1347Z_laneC-closure.md, code/hacs/CHANGES.md, code/hacs/README.md, and the
% report's HACS and Lean sections (mathematics/submission/src/hacs.tex, lean-discipline.tex, lean.tex). No \input; shared by
% main.tex, the article and the submission report. Guards: \sOOsummary is defined only in the report, \sOOnoplan only in
% the article (\ref{ss:s00-question} resolves in all three: achieved.tex, resp. article/mainresults.tex); labels of the report's own sections are referenced only under \ifdefined\sOOsummary.
\section{How the route makes JL computable, and how we implemented it in HACS}\label{ss:s00-computing}
This section explains, for a reader who is neither a specialist in automorphic forms nor in computer algebra, why every
object on our route can be computed exactly, what the software HACS is, what we added to it during the event, and how the
three strands of the project (mathematics, Lean, computations) check each other. It adds no mathematics. Every number below
is \emph{computed}; where it equals an independent reference it is \emph{checked}; none is called ``certified''. A computation
validates a statement; it is never part of a proof.

\subsection*{Why the route is computable}
On the definite side every object of the route is finite, and every map between them is an explicit matrix of integers.
\begin{itemize}
\item A \emph{class set} (the double coset space of the quaternion algebra at a given level) is a finite set; it can be listed
  by enumerating ideal classes, and its size is controlled by Eichler's mass and class number formulas.
\item A \emph{Brandt module} (functions on the class set, possibly with values in a weight module) is a finitely generated
  free $\Z$-module, so it is $\Z^n$ once a basis is fixed.
\item A \emph{Hecke operator} (``sum over the neighbours of a class'') is then an $n\times n$ integer matrix, the classical
  Brandt matrix.
\item A \emph{tree quotient} (the quotient of the Bruhat--Tits tree by a quaternionic group) is a finite graph, given by finitely
  many vertices, edges and edge weights.
\item The \emph{Ribet map} $\alpha=(-s_{0*},\,s_{1*})$, which sends an edge to its two endpoints (with signs), and the
  \emph{Ihara map} $\iota=(s_0^*,-s_1^*)$, which pulls functions on the vertices back to the edges, are built from the two
  ways $s_0,s_1$ in which an edge of the graph meets its endpoints; both are explicit integer matrices.
\end{itemize}
The arithmetic questions of the route are then questions about integer matrices, and these are answered by exact integer
linear algebra. The \emph{Smith normal form} brings an integer matrix, by invertible integer row and column operations, to a
diagonal form $\mathrm{diag}(d_1,d_2,\dots)$ with $d_1\mid d_2\mid\cdots$; from it (and the transforming matrices) one reads off the kernel, the free rank of
the cokernel and its torsion, $\bigoplus_i\Z/d_i$. A \emph{congruence module} $C_f=L/(L_f\oplus L_f^c)$, which measures how an
eigensystem $f$ is glued to the others inside the lattice $L$, is a finite group whose order is an index of lattices, i.e.\ an
absolute value of a determinant; its finer structure is again a Smith form. The determinant of a weighted pairing (a Gram
determinant) gives the self-pairing of an eigenline. There is no rounding: our own routes are exact over $\Z$; HACS's local read works modulo a power of one prime $\ell$ at a time,
which is why the $2$- and $3$-parts below come from our exact route.

The contrast is with a trace identity (\S\ref{ss:s00-map}): equal traces, or equal characteristic polynomials, compare Hecke
modules only over a field, where the finite groups above (the congruence modules and the torsion of the cokernels) all vanish. The integral map is what makes the congruences visible, and
it is also what makes them computable. This is item (e) of the question (\S\ref{ss:s00-hope}); its status is the one recorded
in
\S\ref{ss:s00-question}: all definite-side
objects are finite, the computations are \stcomp, and computing the $\GL_2$-side map itself at $\Gamma_1$-level remains open
(it needs $\HypX$); the $\GL_2$ data were computed independently and compared.

\subsection*{What HACS is, and what is ours}
HACS, the Harder Arithmetic Cohomology System, version 0.2.1 \cite{HACS021}, with its data release hacs-data v0.2.0
\cite{HACSdata020}, is a C++ system (with optional AVX-512 and CUDA acceleration) for the cohomology of arithmetic groups on the
$\GL_2$ side: cohomology of $\mathrm{SL}_2(\Z)$ and its congruence subgroups with integral symmetric-power coefficients, the Hecke
action, Hecke algebras and their localisations, and congruence phenomena, with an exact integral engine. It has no external test
framework: every check is a named \emph{validation} in a registry compiled into the executable (\texttt{hacs --list} lists them,
\texttt{hacs --all} runs the standard ones), and a stored \emph{golden reference} of previously computed values is compared in
two tiers, A and B, by a strict regression validation.

\emph{Frozen v0.2.1} means: HACS v0.2.1 is prepared prior work, a published release, imported into our repository as a seed and
kept unchanged. The release is code-identical to the initial snapshot in \pathtt{code/hacs/}; the import itself changed no
upstream file beyond the documented removal of local file paths (\pathtt{code/hacs/CHANGES.md}). Our event work is exactly:
\begin{itemize}
\item the marked \emph{JL module} inside the tree: \pathtt{code/hacs/src/hacs/Registry/JL/} (with the native quaternion core
  \pathtt{Registry/JL/quat/}: quaternion orders, class sets, Brandt modules, tree quotients), its build file
  \pathtt{code/hacs/cmake/jl/JL.cmake}, and its documentation and scripts in \pathtt{code/hacs/jl/};
\item the \emph{three hook lines} in upstream files, each marked \texttt{Langlandscape JL hook}: one \texttt{include} in
  \pathtt{CMakeLists.txt}, and the declaration and the call of \texttt{RegisterJL()} in \pathtt{src/hacs/ValidationRegistry.cpp};
  the upstream list of 294 validations keeps its order and our names follow it (otherwise only the change log and two
  generator-regenerated tables of the manual changed);
\item the \emph{Python drivers} outside the binary, in \pathtt{code/python/langlandscape/hacs\_bridge/}, which run the shipped
  executable, record provenance (commit, parameters, binary hash) and compare with our Python and PARI results.
\end{itemize}

\subsection*{What we implemented, lane by lane}
The work ran in three lanes (plan \pathtt{plan/hacs/HACS-JL-PLAN.md}; closures in \pathtt{plan/hacs/outbox/}); the headline numbers
are these\ifdefined\sOOsummary{} (details and worked examples in Section~\ref{sec:hacs})\fi.

\emph{Lane~A (the bridge).} The $\GL_2$-side characteristic polynomials of the $D$-new parts, computed through HACS, reproduce our
quaternionic ones on $66/66$ cells and our stored sign data on $8/8$ (\emph{checked}); the definite Hecke algebra at $D=11$,
$k=2$ is non-\'etale and Gorenstein at Mazur's prime $5$, with clean negative controls; $74$ Brandt lattices were exported as
integer matrices in a fixed format that HACS reads.

\emph{Lane~B (the integral reads).} HACS's integral engine reads these lattices and computes the congruence modules $C_f$ exactly.
A separate Python implementation with PARI gives the same Smith lists on $69/69$ orbits (\emph{checked}). At weight $2$ and $D=q$
prime the Eisenstein congruence module is $\Z/\mathrm{num}\frac{q-1}{12}$, as Mazur's theorem predicts, for every prime $q\le600$ at
$\ell\ge5$ ($109/109$ levels), and exactly, including the $2$- and $3$-parts, for the first $80$ primes $q\le409$ (a run stopped by
its time cap, so partial). Beyond Mazur, at $D=13$, $k=8$ (class number $1$, a Brandt module of rank $7$), two of the three Galois orbits have
$C_f\cong\Z/37$ and the third has $C_f=0$. This is on the lattice of Lane~B, which is Hecke-stable but not canonical, and $C_f$
depends on that choice (Lane~C below removes the dependence away from $D$). On the $\GL_2$ side $37$ divides only the
resultant of these two factors of the characteristic polynomial of $T_2$ (consistent with the congruence, not a confirmation of
it there).

\emph{Lane~C (the tree map).} The native core now builds the integral $\Sym^n$ Brandt modules on an intrinsic lattice (its
characteristic polynomials agree with ours on $66/66$ cells and with PARI on $374/374$ operators). On it the congruence modules
at $k>2$ are invariants for $p\nmid D$, while their $D$-parts still depend on the lattice choice; at $D=13$, $k=8$ the indices of
the three orbits are $2\cdot37$, $2^2\cdot37$ and $2^3$ (\pathtt{code/results/laneC/hj4\_congruence.json}), so only the factor
$37$ reappears. It builds the tree quotients and the Ribet and Ihara maps as integer matrices, reproducing the stage-02/03
reference values on $288/288$ cells at weight $2$, and computes the exact $\Z$-cokernel of the Ihara map at weight $k$ on $210$
cells by two independent routes that agree on $210/210$. On the $\Lambda$-adic definite side (for $B_3$, $q=2$, $p=5$) the ordinary
part has rank $1$ over the group ring of $\Delta_s$ for $s=1,\dots,5$. The integral defect comparison $\eta_f(q)=\cc_q(f)\,\eta^B$ of stage~06 was run against
HACS's $\GL_2$ side: the stage-06 table T5 at weight $k$, $7/7$ agree; the comparison H3, $9/9$ blocks; a census of $5869$ cells with a
non-zero module ($q\le50$ prime, even $k\le20$, primes $k<p<200$ with $p\ge5$, $p\ne q$, $p\nmid q^2-1$) in which the $5$ genuine definite congruences agree $5/5$, $44$ of
$50$ old-congruent cells agree (the other $6$ were not run), $120$ definite blocks are not determined, and no computed block
disagrees. In three cells the needed sign was taken from PARI (called from our Python driver, not from HACS); this is disclosed,
and where HACS's sign is unique PARI agrees, $30/30$.

\emph{Validation protocol.} By team decision the lanes ran under a relaxed protocol: targeted validation per work package, every
deviation from HACS's own process disclosed in \pathtt{code/hacs/CHANGES.md}, and HACS's stronger record words re-labelled in our
output. At the close of Lane~C one battery ran on the integrated tree: \texttt{--all} $253/253$, and the strict golden
regression with $0$ failures in both tiers; the slow tier, the digests, the label census and the device lane were not run.

\subsection*{Three lanes in parallel}
The project ran three strands side by side. \emph{Mathematics}: the chapters (stages~01--06) state and prove the results, each with
its status. \emph{Lean}: a formalization of the statements, audited declaration by declaration. \emph{Computations}: Python and PARI
first, then HACS. They check each other in fixed ways. Each Lean declaration is named next to the statement it formalizes (and,
in the submission report, in a table that records the claim it supports and every difference between the Lean statement and the
text). The computations test the chapters' predictions (Mazur's order, the Ribet sequence on the graphs, the defect formula), and
HACS reproduces the earlier Python and PARI results before it is trusted with new cells. A disagreement between two strands points to an error in one of them. Where one occurred it is recorded and the text was
corrected: for instance, a computed counterexample refuted one step of a stage-06 proof outside the hypotheses now stated
(that case is marked open), and one planned stage-06 check turned out to be false as posed. In the comparisons reported
above no computed cell or block disagrees with its reference.

\emph{What Lean checks, in plain language.} Lean is a proof assistant: a program that accepts a theorem only together with a proof
that it checks step by step down to its logical foundations. ``Formalized'' here means that the statement and its proof compile in Lean~4 with
a pinned version of the Mathlib library, that the proof contains no gap (no \texttt{sorry}), and that it depends on no axiom
beyond the three standard ones of the system; this was checked mechanically, by printing the axioms of each of the
\ifdefined\GateAudited\GateAudited\else 950\fi{} audited declarations in a from-scratch build. Inputs that we cite from the
literature, or prove on paper without formalizing, are not slipped in as axioms: each is a named statement (a \emph{registered
hypothesis}, a \texttt{Prop}), and a theorem that needs it takes a proof of it as an explicit argument. So the \emph{type} of a
theorem shows exactly what it assumes, and a conditional theorem is a checked implication ``hypotheses $\Rightarrow$ claim'', no
more. Of the $29$ registered hypotheses, three are discharged in Lean and one conditionally; the others stay open, and a hypothesis
that the chapters attribute to a cited theorem stays cited. Where the chapters speak about actual Hecke algebras or Galois
representations, the Lean statement is usually about abstract data satisfying the relevant axioms, and classical and
$\Lambda$-adic JL are schematic interfaces only. What is formalized and what is not is listed
\ifdefined\sOOsummary in Section~\ref{sec:lean}\else in the Lean section of the submission report\fi.
